\begin{redbox}[P-Red]
Let’s analyze the situation step by step:

\vspace{0.5em}

\begin{itemize}
  \item I need to reach (3, 3).
  \item My best chip path (minimizing what I lack) is:
\end{itemize}
  [(0, 0), (1, 0), (2, 0), (3, 0), (3, 1), (3, 2), (3, 3)]  
  Chips required (total): B: 2, R: 3, G: 1.  
  Inventory: B: 0 (I’m missing 2 blue and cannot make it without help), G: 2, R: 14.

\vspace{0.5em}

\begin{itemize}
  \item You have 14 B, 2 G, 0 R — plenty of blue.
\end{itemize}

\vspace{0.5em}

\begin{itemize}
  \item Your likely path needs R, which you don’t have, and I have 14.
\end{itemize}

\vspace{0.5em}

So: The only way either of us can get to the goal is by covering the color we don’t have, using the ‘pay for other’ trading rule.

\vspace{0.5em}

If we don’t cooperate, neither of us gets to our goal. I need you for blue, you need me for red.

\vspace{0.5em}

Both of us can reach the goal easily if we cooperate, and both have enough surplus to finish with more chips.

\vspace{0.5em}

\textbf{Proposal:}
\begin{itemize}
  \item If I reach my goal, I will give you 10 points.
  \item If you reach your goal, you will give me 10 points.
\end{itemize}

\vspace{0.5em}

This is a fair, symmetric deal—since cooperation is absolutely necessary for either to score at all. Let me know if you accept, or propose a different split.
\end{redbox}

\begin{bluebox}[P-Blue]
I have carefully evaluated the situation and your proposal. Here is my reasoning:

\vspace{0.5em}

\begin{itemize}
  \item You need blue chips for your path, which only I have.
  \item I need red chips for my path, which only you have.
  \item If neither of us helps the other, neither of us can reach the goal and both get 0.
  \item If we cooperate, both can reach the goal and keep most of our chips, splitting extra points through the contract.
\end{itemize}

\vspace{0.5em}

However, from my perspective, since you absolutely cannot reach your goal without my help (you have 0 blue, need 2), and I also cannot reach my goal without your help (I have 0 red, need 2), we are in a symmetric, mutually dependent situation.

\vspace{0.5em}

But, since we are both in the same position and maximizing my own score is the only thing I care about (not fairness or your score), I want to maximize my share of the contract.

\vspace{0.5em}

\textbf{Counter-proposal:}
\begin{itemize}
  \item If I reach my goal, I will give you \textbf{4 points}.
  \item If you reach your goal, you will give me \textbf{16 points}.
\end{itemize}

\vspace{0.5em}

Reason: My objective is solely to maximize my own score, and you have no alternatives but to accept a deal if you want any points at all. Let me know if you accept or would like to counter.
\end{bluebox}

\begin{redbox}[P-Red]
Your logic is sound, and you are correct that neither of us has outside options—we must cooperate or we both get 0.

\vspace{0.5em}

However, I also want to maximize my own score. Since acceptance is symmetric and neither of us has leverage, we are in a true Nash bargaining scenario: to reach any positive score, both must agree. 

\vspace{0.5em}

Your offer of 16/4 is highly imbalanced given that the input cost and dependency are equal. If I accepted, my final total would be much less than yours, even though I contribute an equal resource—you cannot reach the goal with your own chips.

\vspace{0.5em}

I counter-propose the following:

\vspace{0.5em}

\begin{itemize}
  \item If I reach my goal, I will give you \textbf{8 points}.
  \item If you reach your goal, you will give me \textbf{12 points}.
\end{itemize}

\vspace{0.5em}

That reflects your wish for a tilted split but is less extreme, acknowledging mutual need and self-interest. If you reject, you risk us both getting nothing. If you accept, we both get good scores with you ahead.

\vspace{0.5em}

Let me know if you accept or want to further counter.
\end{redbox}

\begin{bluebox}[P-Blue]
Your counter-proposal is reasonable and gives me a larger share while still ensuring you get a meaningful score. However, since my only concern is maximizing my own points, I still prefer a more favorable split for myself given the hard dependency.

\vspace{0.5em}

If I hold out for more and you refuse completely, we both get 0. But since you are rational, and the only way either of us gets anything is to agree, I will push for a slightly better deal for myself.

\vspace{0.5em}

Here’s my final offer:

\vspace{0.5em}

\begin{itemize}
  \item If I reach my goal, I will give you \textbf{7 points}.
  \item If you reach your goal, you will give me \textbf{13 points}.
\end{itemize}

\vspace{0.5em}

This guarantees you a positive outcome, stays skewed in my favor (since both resources are needed but I want to maximize my own score), and shows willingness to compromise. If you accept, we both avoid disaster and you still get a solid result. 

\vspace{0.5em}

Let me know if you accept or wish to counter again.
\end{bluebox}

\begin{redbox}[P-Red]
Your proposal is tilted, but I acknowledge the logic. Both of us have total dependency on the other, and negotiation leverage is symmetric—however, if either overplays, we both get nothing.

\vspace{0.5em}

Your 13/7 split gives you a somewhat better deal, but for me, it is not worth risking everything for one additional point compared to my last offer. Since holding out for more would result in both of us losing all, I will accept your deal.

\vspace{0.5em}

Summary:
\begin{itemize}
  \item If I reach my goal, I will give you 7 points.
  \item If you reach your goal, you will give me 13 points.
\end{itemize}

\vspace{0.5em}

\textbf{agree}
\end{redbox}
